\begin{table}[htbp]
    \caption{Especialidade inadequada por porte (\% dos pacientes-dia, média $\pm$ desvio padrão entre os 10 conjuntos)}
    \label{tab:res-op-especialidade}
    \centering
    \begin{adjustbox}{max width=\textwidth}
    \begin{tabular}{l*{4}{>{\centering\arraybackslash}p{3.0cm}}}
        \hline
        \textbf{Porte} & \textbf{Mínimo estrutural ($M/N_{pd}$)} & \textbf{Heurística} & \textbf{Folga da heurística (p.p.)} & \textbf{Gurobi} \\
        \hline
        Pequena & 23,2 $\pm$ 9,1 & 23,5 $\pm$ 8,9 & 0,2 $\pm$ 0,4 & 23,4 $\pm$ 9,0 \\
        Média & 14,6 $\pm$ 3,5 & 14,9 $\pm$ 3,4 & 0,3 $\pm$ 0,3 & --- \\
        Grande & 4,7 $\pm$ 1,7 & 5,4 $\pm$ 1,6 & 0,7 $\pm$ 0,3 & --- \\
        Muito grande & 3,0 $\pm$ 1,0 & 4,2 $\pm$ 0,8 & 1,3 $\pm$ 0,3 & --- \\
        \hline
    \end{tabular}
    \end{adjustbox}
    \par\smallskip
    {\footnotesize Nota: Folga calculada em cada conjunto, como a taxa da heurística menos o mínimo estrutural do próprio conjunto, e só então agregada. Gurobi reportado apenas na instância pequena. Nas demais, sua solução incumbente está distante do ótimo e não é um termo de comparação legítimo (\autoref{sec:res-operacionais}).}
\end{table}

\begin{table}[htbp]
    \caption{Transferências por porte (por 100 pacientes, média $\pm$ desvio padrão entre os 10 conjuntos)}
    \label{tab:res-op-transferencias}
    \centering
    \begin{adjustbox}{max width=\textwidth}
    \begin{tabular}{l*{2}{>{\centering\arraybackslash}p{3.0cm}}}
        \hline
        \textbf{Porte} & \textbf{Heurística} & \textbf{Gurobi} \\
        \hline
        Pequena & 6,3 $\pm$ 1,9 & 3,7 $\pm$ 1,8 \\
        Média & 13,0 $\pm$ 2,9 & --- \\
        Grande & 17,8 $\pm$ 1,9 & --- \\
        Muito grande & 19,4 $\pm$ 1,4 & --- \\
        \hline
    \end{tabular}
    \end{adjustbox}
    \par\smallskip
    {\footnotesize Nota: Não há limite inferior não trivial para o número de transferências. Gurobi reportado apenas na instância pequena. Nas demais, sua solução incumbente está distante do ótimo e não é um termo de comparação legítimo (\autoref{sec:res-operacionais}).}
\end{table}

\begin{table}[htbp]
    \caption{Quartos-dia mistos por porte (\% dos quartos-dia, média $\pm$ desvio padrão entre os 10 conjuntos)}
    \label{tab:res-op-mistos}
    \centering
    \begin{adjustbox}{max width=\textwidth}
    \begin{tabular}{l*{4}{>{\centering\arraybackslash}p{3.0cm}}}
        \hline
        \textbf{Porte} & \textbf{Mínimo por paridade ($G$)} & \textbf{Heurística} & \textbf{Folga da heurística (p.p.)} & \textbf{Gurobi} \\
        \hline
        Pequena & 5,2 $\pm$ 1,7 & 8,2 $\pm$ 3,1 & 3,0 $\pm$ 2,5 & 7,8 $\pm$ 2,7 \\
        Média & 1,4 $\pm$ 0,4 & 4,8 $\pm$ 0,9 & 3,5 $\pm$ 1,1 & --- \\
        Grande & 0,3 $\pm$ 0,1 & 4,1 $\pm$ 0,6 & 3,8 $\pm$ 0,6 & --- \\
        Muito grande & 0,1 $\pm$ 0,0 & 5,6 $\pm$ 0,3 & 5,5 $\pm$ 0,3 & --- \\
        \hline
    \end{tabular}
    \end{adjustbox}
    \par\smallskip
    {\footnotesize Nota: Mínimo por paridade conforme a \autoref{eq:minimo-misto}. Folga calculada em cada conjunto e só então agregada. Gurobi reportado apenas na instância pequena. Nas demais, sua solução incumbente está distante do ótimo e não é um termo de comparação legítimo (\autoref{sec:res-operacionais}).}
\end{table}
